\documentclass[11pt]{article}
\usepackage{booktabs}
\usepackage{amsmath}
\usepackage{natbib}
\usepackage[margin=1in]{geometry}

\title{Three-Line Tables and Honest Baselines}
\author{Example Author}
\date{}

\begin{document}
\maketitle

\begin{abstract}
A short clean manuscript used to show what a passing build looks like.
\end{abstract}

\section{Introduction}
\label{sec:intro}
Fusion readers improve retrieval-augmented question answering
\citep{lewis2020rag}. Dense retrieval is a strong baseline
\citep{karpukhin2020dpr}. We compare both under one budget.

\section{Method}
\label{sec:method}
Given a question $q$ and the top-$k$ passages, the reader scores each passage
and attends over the scores. Training uses the frozen retriever.

\section{Results}
\label{sec:results}
Table~\ref{tab:dev} reports development accuracy. The gap between the sparse
and dense baselines is 9.8 points, consistent with prior work.

\begin{table}[t]
\centering
\caption{Development accuracy on the synthetic split.}
\label{tab:dev}
\begin{tabular}{lcc}
\toprule
Model & Dev & Test \\
\midrule
Sparse only & 61.2 & 60.4 \\
Dense only & 71.0 & 70.2 \\
Fusion (ours) & 72.4 & 71.8 \\
\bottomrule
\end{tabular}
\end{table}

Section~\ref{sec:method} fixes the training budget, so the comparison in
Table~\ref{tab:dev} is like for like.

\bibliographystyle{plainnat}
\bibliography{refs}

\end{document}
